% !PRINTED 607
\weberpage{607}
%<<CONT
durch $p$ theilbar sein, und dies ist auch die hinreichende Bedingung.

2. \gesperrt{Es ist also $p$ nur dann mit dem Quadrat eines Primfunctionals associirt, wenn $p$ in der Grundzahl $\Delta$ aufgeht, in Uebereinstimmung mit dem allgemeinen Satze (\S.\,158).}

\begin{center}
\S.\,166.

\gesperrt{Functionale im quadratischen K\"orper und quadratische Irrationalzahlen.}
\end{center}

Wir betrachten nun in unserem quadratischen K\"orper ein beliebiges ganzes Functional und bilden seine Basisform
\[
\varphi = \alpha_1 t_1 + \alpha_2 t_2,
\]
indem wir $\alpha_1, \alpha_2$ nach \S.\,165 bestimmen:
\weq{1}{\alpha_1 = a_{1,1}, \quad \alpha_2 = a_{1,2} + a_{2,2}\Theta.}
Nach \S.\,146 ist dann
\weq{2}{N_a(\varphi) = a_{1,1}a_{2,2}.}
Geben wir der Form $\varphi$ den Ausdruck
\weq{3}{\varphi = a_{1,1}t_1 + (a_{1,2}+a_{2,2}\Theta)t_2,}
und bilden die Norm, so ergiebt sich
\weq{4}{N(\varphi) = a_{1,1}^2 t_1^2 + [2a_{1,1}a_{1,2} + a_{1,1}a_{2,2}(\Theta+\Theta')]t_1 t_2 + t_2^2 N(\alpha_2)}
und der Theiler dieser Form ist also nach (2) gleich $a_{1,1}a_{2,2}$.
Wenn wir also
\weq{5}{\begin{aligned}
a_{1,1}^2 &= c\,a_{1,1}a_{2,2}\\
2a_{1,1}a_{1,2}+a_{1,1}a_{2,2}(\Theta+\Theta') &= b\,a_{1,1}a_{2,2}\\
N(\alpha_2) &= -a\,a_{1,1}a_{2,2}
\end{aligned}}
setzen, so sind $a, b, c$ ganze rationale Zahlen ohne gemeinschaftlichen Theiler, und der Quotient
\weq{6}{\frac{\alpha_2}{\alpha_1} = \omega}
ist eine Wurzel der quadratischen Gleichung
\weq{7}{c\omega^2 = a + b\omega.}
% [sic: The source prints the incomplete subscript $a_{1,}$ here; equations (1) and (5) identify the intended quantity as $a_{1,2}$.]
Wenn wir in den Ausdr\"ucken (1) die Werthe von $a_{1,1}$ und $a_{1,}$ aus (5) substituiren, so folgt
\weq{8}{\begin{aligned}
2\alpha_2 &= a_{2,2}(b+\Theta-\Theta')\\
\alpha_1 &= c\,a_{2,2}.
\end{aligned}}
